\documentclass[12pt]{article}
\usepackage[margin=1in]{geometry}
\usepackage{enumitem}
\usepackage{titlesec}

\titleformat{\section}{\large\bfseries}{}{0em}{}
\titlespacing*{\section}{0pt}{1em}{0.3em}

\setlength{\parindent}{0pt}
\setlength{\parskip}{0.5em}

\begin{document}

\begin{center}
    {\LARGE \textbf{IcePEAK}}\\[0.3em]
    {\large CIS 5680 --- Group Game Project 2}
\end{center}

\section{High Concept}

Conquer a treacherous frozen mountain in VR by driving ice picks into glacial walls with your own hands, managing dwindling stamina and supplies as you race against avalanches and crumbling ice to reach the summit alive.

\section{Genre}

First-person VR survival-adventure with resource management and physics-based climbing.

\section{Features}

\begin{itemize}[leftmargin=2em]
    \item \textbf{Physical VR climbing.} Players use motion controllers to swing and plant ice picks into the mountainside. Each grip, reach, and pull feels deliberate and physical --- you climb with your own arms.

    \item \textbf{Destructible ice surfaces.} Most iced surfaces crack and shatter roughly two seconds after an ice pick is driven in, forcing players to move quickly and plan their next hold before the current one gives way.

    \item \textbf{Unbreakable rock platforms.} Scattered stone ledges and rocky outcrops serve as safe resting points where players can catch their breath, check inventory, and plan the next stretch of the ascent.

    \item \textbf{Resource management.} A stamina gauge depletes with every swing and hold. Hunger and body heat drain over time. Players scavenge supplies like rations and heat packs from caches embedded in the cliffside.

    \item \textbf{Dynamic environmental hazards.} Falling rocks, sudden ice storms, and cracking ice shelves keep the climber on edge. Audio and visual cues warn of incoming danger, rewarding alert players who reposition in time.

    \item \textbf{Route-finding strategy.} Multiple paths lead upward, each with different risk--reward trade-offs: a direct overhang is fast but stamina-draining, while a longer traverse passes more supply caches but exposes the player to rockfall.

    \item \textbf{Immersive body presence.} Gear and collected items are visible on the player's body --- carabiners on the harness, a backpack with supplies, ice picks in hand --- reinforcing the feeling of being on the mountain.

    \item \textbf{Escalating altitude.} As the player climbs higher, conditions worsen: visibility drops, wind intensifies, and ice becomes more fragile, creating a natural difficulty curve tied to elevation.
\end{itemize}

\section{Player Challenge / Motivation}

The player is a solo mountaineer attempting to summit a massive frozen peak. The challenge is both physical and strategic: every handhold is a split-second decision, and every resource spent is one fewer for the dangers ahead. The drive to see what lies above the next ridge --- and the ever-present threat of a fatal fall --- keeps the player pushing upward.

\section{Design Goals}

\textbf{Tension and suspense.} The two-second ice timer and limited stamina create constant pressure. Players are never fully safe, and each stretch of climbing carries real stakes --- a missed grab or a crumbling hold means falling to the last safe platform.

\textbf{Physical immersion.} VR hand tracking makes climbing feel earned. The act of reaching overhead, swinging a pick, and pulling yourself up engages the player's body in a way flat-screen games cannot replicate.

\textbf{Strategic depth.} Route selection, stamina budgeting, and supply management give the game a planning layer beneath the action. Experienced players will learn to read the mountain and find optimal paths.

\section{Unique Selling Points}

\begin{itemize}[leftmargin=2em]
    \item \textbf{Destructible-ice climbing} is a mechanic not seen in existing VR titles. The ticking clock on every handhold transforms climbing from a relaxed traversal into a tense, rhythmic puzzle.
    \item \textbf{Full body-mounted inventory in VR} lets players physically grab supplies off their own harness, deepening immersion beyond typical menu-based systems.
    \item \textbf{Elevation-driven difficulty} ties challenge progression to the physical act of climbing rather than artificial level gates, making the mountain itself the game's narrative arc.
\end{itemize}

\end{document}
